\documentclass[11pt]{article}
\usepackage[margin=1in]{geometry}
\usepackage{longtable}
\usepackage{array}
\usepackage{amsmath}
\setlength{\parindent}{0pt}
\setlength{\parskip}{0.6em}

\begin{document}

\section*{Review of \texttt{main.tex}: SENA manuscript}

\textbf{Overall assessment.}
The manuscript is strong in scope and technical ambition, with a coherent probabilistic structure and many implementation-oriented details.
The highest-risk issues are not arithmetic mistakes, but specification ambiguities that can produce materially different implementations (especially interval indexing, state initialization, and symbol overload).
I did not find clear algebraic sign/factor errors in the core inventory, lead-time mapping, or safety-stock equations.

\section*{Most important technical findings}

\begin{enumerate}
  \item \textbf{Definite error: off-by-one/index-domain inconsistency in the age-state recursion.}

  \textbf{Location:} Section ``Restocking as a policy with an explicit order pipeline,'' age update equation for $\text{age}_{k+1,i}$; compare with earlier definition that transition-event quantities (including $G_{k,i}$) are defined for $k=2,\dots,K$ only.

  \textbf{Problem:} The manuscript defines transition events on $k=2,\dots,K$, but writes
  \[
    \text{age}_{k+1,i}=\begin{cases}
      0,& G_{k,i}=1,\\
      \text{age}_{k,i}+\Delta t_k,& G_{k,i}=0.
    \end{cases}
  \]
  Without an explicit $k$ range, this is inconsistent either at $k=1$ (requires undefined $G_{1,i}$) or at $k=K$ (creates unused $\text{age}_{K+1,i}$ while leaving $\text{age}_{2,i}$ underspecified).

  \textbf{Why it matters:} This directly affects placement propensity $\pi^{\mathrm{ord}}_{k,i}$ and therefore order timing and inferred stockout risk.

  \textbf{Suggested fix:} Re-index explicitly, for example
  \[
    \text{age}_{k,i}=\begin{cases}
      0,& G_{k-1,i}=1,\\
      \text{age}_{k-1,i}+\Delta t_{k-1},& G_{k-1,i}=0,
    \end{cases}\quad k=2,\dots,K,
  \]
  with $\text{age}_{1,i}=a_i^{\mathrm{init}}$.

  \item \textbf{Likely issue: missing initial-state priors for multiple random-walk latent processes.}

  \textbf{Location:} Drift and lead-time state equations for $\log\lambda_{k,j}$, $\log\mu_{k,i}$, $\mu^{\mathrm{use}}_{k,j,i}$, $\ell_{k,i}$, and $\psi_{k,i}$.

  \textbf{Problem:} Recursions are specified, but explicit priors/initial conditions at $k=1$ are not specified for these states (in contrast to $x_{1,i}$, $B_{1,i}$, and $\text{age}_{1,i}$, which are initialized explicitly).

  \textbf{Why it matters:} Different implicit defaults (fixed constants, diffuse priors, empirical Bayes starts) can materially change posterior behavior and reproducibility.

  \textbf{Suggested fix:} Add a compact ``initial latent-state priors'' block, e.g.
  $\log\lambda_{1,j}\sim\text{Normal}(m_{\lambda,j},s_{\lambda,j}^2)$,
  $\log\mu_{1,i}\sim\text{Normal}(m_{\mu,i},s_{\mu,i}^2)$,
  $\ell_{1,i}\sim\text{Normal}(m_{L,i},s_{L,i}^2)$,
  $\psi_{1,i}\sim\text{Normal}(m_{v,i},s_{v,i}^2)$,
  and (if used) $\mu^{\mathrm{use}}_{1,j,i}\sim\text{Normal}(m_{\mathrm{use},j,i},s_{\mathrm{use},j,i}^2)$.

  \item \textbf{Likely issue: lag-index boundary ambiguity in delayed lead-time observations.}

  \textbf{Location:} Optional lead-time observations with
  $\ell_{k-\delta^m_{k,i},i}$ and $\psi_{k-\delta^c_{k,i},i}$, where $\delta\in\{0,\dots,D\}$.

  \textbf{Problem:} The text does not constrain $\delta$ relative to $k$, so early intervals can reference index $\le 0$.

  \textbf{Why it matters:} This creates implementation divergence (padding conventions vs truncation) and can silently bias early-interval likelihood terms.

  \textbf{Suggested fix:} State one explicit boundary rule, e.g. enforce $\delta\le k-1$, or map to
  $\ell_{\max(1,k-\delta),i}$ and $\psi_{\max(1,k-\delta),i}$.

  \item \textbf{Notation/terminology drift check (required): conflict in demand vs realized-consumption naming.}

  \textbf{Location:} Inventory dynamics and output sections.

  \textbf{Problem:} Symbols with the $C$ prefix are used for conceptually different roles:
  $C_{k,i}$ is realized/capped consumption, while $C^{\mathrm{svc}}_{k,i}$ and $C^{\mathrm{ret}}_{k,i}$ are used inside
  $D_{k,i}=C^{\mathrm{svc}}_{k,i}+C^{\mathrm{ret}}_{k,i}$ and later described as latent demand-side components.

  \textbf{Why it matters:} This can cause implementation bugs in censoring, KPI computation, and interpretation of ``consumption'' outputs.

  \textbf{Suggested fix:} Reserve $C$ for realized quantities only; rename latent demand-side components to
  $D^{\mathrm{svc}}_{k,i}$ and $D^{\mathrm{ret}}_{k,i}$, and keep
  $C_{k,i}=\min\{D_{k,i},x_{k-1,i}+R_{k,i}\}$ as the realized counterpart.

  \item \textbf{Likely issue: symbol overload for $B$ creates implementation-facing ambiguity.}

  \textbf{Location:} Global indices/definitions and recipe subsection.

  \textbf{Problem:} $B$ denotes number of bundle types, $B_{k,i}$ denotes in-transit inventory state, and $B^{(m)}_{k,j,i}$ denotes per-instance Bernoulli usage indicator.

  \textbf{Why it matters:} Overloaded symbols increase risk of coding errors and documentation drift, especially in appendices or software implementations.

  \textbf{Suggested fix:} Rename at least one family, e.g. bundle count $N_{\mathrm{bundle}}$, indicator $Z^{(m)}_{k,j,i}$, while keeping $B_{k,i}$ for pipeline stock.

  \item \textbf{Unsupported-claim calibration issue in promo/bundle narrative.}

  \textbf{Location:} Promotions and bundles motivation/conclusion statements (e.g., claims that explicit promo modeling ``can improve identifiability'' and ``improving both inference stability and business interpretability").

  \textbf{Problem:} These are plausible claims, but currently stronger than the evidence shown in this manuscript draft (no quantitative ablation results are presented in the same section).

  \textbf{Why it matters:} Overstated claims weaken technical credibility at review time.

  \textbf{Suggested fix:} Keep causal language conditional (``can''/``hypothesized to'') and explicitly point to the Evaluation protocol for publication readiness section metrics as the validation mechanism.
\end{enumerate}

\section*{Symbol drift snapshot (requested check)}

\begin{longtable}{|>{\raggedright\arraybackslash}p{0.14\textwidth}|>{\raggedright\arraybackslash}p{0.26\textwidth}|>{\raggedright\arraybackslash}p{0.28\textwidth}|>{\raggedright\arraybackslash}p{0.24\textwidth}|}
\hline
\textbf{Symbol} & \textbf{First definition / descriptor} & \textbf{Later descriptor variant} & \textbf{Assessment} \\
\hline
$C_{k,i}$ & Realized consumption/usage capped by available stock. & Coexists with $C^{\mathrm{svc}}_{k,i},C^{\mathrm{ret}}_{k,i}$ that are later used to build latent demand $D_{k,i}$. & Conflict risk (realized vs latent semantics under same prefix). \\
\hline
$B$ / $B_{k,i}$ / $B^{(m)}_{k,j,i}$ & Bundle-count cardinality, in-transit pipeline stock, and Bernoulli usage indicator. & All three appear as central symbols in separate model layers. & Overload likely to confuse implementation and exposition. \\
\hline
$D_{k,i}$ & Latent unconstrained demand/usage for interval $k$. & Later decomposition uses $C^{\mathrm{svc}}+C^{\mathrm{ret}}$ naming. & Descriptor drift via mixed latent/realized naming family. \\
\hline
\end{longtable}

\section*{Meaningful editorial findings}

\begin{enumerate}
  \item \textbf{Location:} Equation-recursion blocks across Sections 2, 4, and 6.

  \textbf{Problem:} Several recursions omit explicit index ranges right after equations, which is the root cause of multiple ambiguities above.

  \textbf{Why it matters:} Readers can mis-implement edge intervals even when core equations are correct.

  \textbf{Suggested fix:} Add one short sentence after each recurrence: ``valid for $k=2,\dots,K$'' (or corresponding range).

  \textbf{Candidate rewrite:} ``For all recursions below, $k=2,\dots,K$ unless otherwise stated.''

  \item \textbf{Location:} Promotions/bundles section concluding sentence.

  \textbf{Problem:} The sentence states improvement outcomes without immediate quantitative support.

  \textbf{Why it matters:} Tone overreach is likely to draw reviewer pushback.

  \textbf{Suggested fix:} Downtone to hypothesis language and reference planned ablations in the Evaluation protocol for publication readiness section.

  \textbf{Candidate rewrite:} ``In this formulation, explicit bundles are expected to improve identifiability and interpretability; the Evaluation protocol for publication readiness section specifies ablations to test this claim.''
\end{enumerate}

\section*{Proofing sweep (mandatory micro-scan)}

I ran the required scanner and spot-checked results:

\begin{itemize}
  \item \textbf{Location:} command-level scan of \texttt{main.pdf} via \texttt{python scripts/proofing\_scan.py main.pdf --max-hits 80}.

  \textbf{Problem:} Hits were almost entirely \texttt{PUNC\_DOUBLE\_PERIOD} from table-of-contents dot leaders (``. . .''), i.e., false positives for this rule.

  \textbf{Suggested fix:} No manuscript text fix needed; if desired, refine scanner to ignore TOC pages.

  \item \textbf{Location:} scan of source text and bibliography via \texttt{python scripts/proofing\_scan.py main.tex --max-hits 80} and \texttt{python scripts/proofing\_scan.py refs.bib --max-hits 80}.

  \textbf{Problem:} No high-confidence pattern-scan hits in manuscript source or bibliography entries.

  \textbf{Suggested fix:} None required for the scanned rule set.
\end{itemize}

\section*{Highest-priority fixes (in order)}

\begin{enumerate}
  \item Fix age-state indexing/domain ambiguity (definite correctness issue).
  \item Add explicit initial priors for all random-walk latent states.
  \item Resolve $C$-family and $B$-family notation overload before implementation appendices/code release.
  \item Calibrate promo/bundle benefit claims to match currently presented evidence.
\end{enumerate}

\section*{Items needing external verification}

No external factual claims were identified that require outside source verification for this review.
The flagged issues are internal consistency, identifiability, indexing, and claim-strength calibration issues within the manuscript itself.

\end{document}